\documentclass[conference]{IEEEtran}
\IEEEoverridecommandlockouts
\usepackage{cite}
\usepackage{amsmath,amssymb,amsfonts}
\usepackage{graphicx}
\usepackage{textcomp}
\usepackage{xcolor}
\usepackage[utf8]{inputenc}

\def\BibTeX{{\rm B\kern-.05em{\sc i\kern-.025em b}\kern-.08em
    T\kern-.1667em\lower.7ex\hbox{E}\kern-.125emX}}

\begin{document}

\title{Sistem Informasi Pemetaan Kapasitas Tempat Penampungan Sampah Sementara (TPS) Berbasis Web di Kota Surabaya\\
    \large \textit{(Draft --- Bab 1 dari 4: Latar Belakang)}}

\author{
    \IEEEauthorblockN{Bayu Wirayudha
        (3125521028) \\
        Muhammad Hanif Wijayanto (3125521029)}
    \IEEEauthorblockA{\textit{Program Studi Teknik Informatika} \\
        \textit{Politeknik Elektronika Negeri Surabaya}\\
    }
}

\maketitle

\begin{abstract}
    % CATATAN: abstrak ini masih draft kasar dan sebaiknya ditulis ulang
    % setelah Bab 2--4 selesai, karena abstrak IEEE idealnya merangkum
    % seluruh isi paper (masalah, metode, hasil), bukan hanya latar belakang.
    Pengelolaan sampah di kota besar seperti Surabaya menghadapi tantangan
    signifikan akibat volume timbulan yang tinggi dan keterbatasan
    transparansi data kapasitas Tempat Penampungan Sampah Sementara (TPS).
    Makalah ini mengusulkan sebuah sistem informasi berbasis web yang
    memetakan sebaran, status, dan kapasitas TPS di Kota Surabaya untuk
    mendukung pengambilan keputusan oleh warga, petugas kebersihan, dan
    perencana kota. \textit{[Ringkasan metode dan hasil akan ditambahkan
            setelah Bab 3 dan Bab 4 selesai.]}
\end{abstract}

\begin{IEEEkeywords}
    sistem informasi geografis, pengelolaan sampah, e-governance, tempat penampungan sementara, smart city
\end{IEEEkeywords}

\section{Pendahuluan}

\subsection{Konteks dan Permasalahan}

Sebagai salah satu kota metropolitan terbesar di Indonesia, Surabaya
menghasilkan volume sampah harian yang signifikan. Dinas Lingkungan
Hidup (DLH) Kota Surabaya secara konsisten mencatat volume sampah yang
masuk ke Tempat Pemrosesan Akhir (TPA) Benowo berada pada kisaran
1.500--1.700 ton per hari dalam kondisi normal, dengan komposisi
didominasi oleh sampah organik sekitar 60\% dari total volume
\cite{antara2023,liputan62023}. Seluruh volume tersebut terlebih
dahulu dikumpulkan di ratusan Tempat Penampungan Sampah Sementara
(TPS) yang tersebar di lebih dari 30 kecamatan sebelum diangkut menuju
TPA. Berdasarkan pemberitaan terbaru, TPA Benowo --- yang telah
beroperasi sejak tahun 2001 dan menampung sampah domestik dari sekitar
190 TPS di Surabaya \cite{kompas2026} --- mengalami insiden kebakaran
pada Juli 2026, yang kemudian memicu sorotan DPRD Kota Surabaya
terhadap keberlanjutan skema pengelolaan sampah kota dan mendorong
rencana transisi menuju teknologi pengolahan yang lebih modern
\cite{suararakyat2026}. Peristiwa ini menegaskan bahwa infrastruktur
persampahan kota saat ini masih rentan terhadap gangguan operasional,
sehingga memperkuat urgensi tersedianya sistem pemantauan dan
pengelolaan data TPS yang lebih transparan dan berbasis data.

Meskipun jaringan TPS memegang peranan krusial sebagai simpul awal
rantai pasok pengelolaan sampah kota, pengelolaan informasi mengenai
sebaran, kapasitas, dan cakupan layanan TPS terhadap wilayah kelurahan
masih dilakukan secara manual dan tersimpan dalam berkas dokumen
statis (spreadsheet) yang terfragmentasi. Kondisi ini menimbulkan
beberapa persoalan operasional utama, yang dirangkum pada Tabel
\ref{tab:masalah}.

\begin{table}[h]
    \caption{Ringkasan Permasalahan Eksisting}
    \label{tab:masalah}
    \centering
    \begin{tabular}{p{1.7cm}p{2.4cm}p{3.4cm}}
        \hline
        \textbf{Domain}      & \textbf{Kondisi Eksisting}                                   & \textbf{Dampak Operasional}                               \\
        \hline
        Pengawasan kapasitas & Pemantauan manual tanpa data kubikasi presisi                & Keterlambatan deteksi TPS yang melebihi kapasitas         \\
        \hline
        Data spasial         & Data lokasi \& atribut TPS terfragmentasi di berkas terpisah & Sulit memetakan zona rawan akumulasi sampah secara akurat \\
        \hline
        Keterbukaan publik   & Warga sulit memperoleh informasi lokasi \& status TPS aktif  & Berpotensi memicu pembuangan sampah di luar lokasi resmi  \\
        \hline
    \end{tabular}
\end{table}

\subsection{Motivasi Solusi Berbasis E-Governance}

Persoalan di atas menunjukkan adanya kesenjangan antara ketersediaan
data operasional persampahan dan aksesibilitasnya bagi pemangku
kepentingan yang membutuhkan --- baik warga yang ingin mengetahui
lokasi TPS terdekat, petugas yang perlu memantau kondisi lapangan,
maupun perencana kota yang memerlukan data agregat untuk pengambilan
keputusan. Pendekatan \textit{e-governance} berbasis Sistem Informasi
Geografis (SIG/GIS) menawarkan solusi untuk mendigitalkan data
persampahan yang sebelumnya statis menjadi informasi yang dapat
divisualisasikan, dicari, dan dianalisis secara interaktif, sejalan
dengan inisiatif portal data terbuka pemerintah seperti Sistem
Informasi Pengelolaan Sampah Nasional (SIPSN) yang dikelola Kementerian
Lingkungan Hidup dan Kehutanan \cite{sipsn}, maupun kajian global
mengenai pengelolaan sampah perkotaan oleh World Bank \cite{worldbank}.

\subsection{Tujuan}

Berdasarkan uraian permasalahan tersebut, makalah ini bertujuan untuk:
\begin{enumerate}
    \item Merancang dan membangun sistem informasi berbasis web yang
          memetakan sebaran, status, dan kapasitas TPS di Kota Surabaya.
    \item Memodelkan relasi TPS terhadap kelurahan yang dilayani,
          mengingat satu TPS dapat melayani lebih dari satu kelurahan
          sekaligus berdasarkan data internal yang diperoleh.
    \item Menyediakan mekanisme validasi data dan pengujian sistem
          untuk memastikan keandalan informasi yang disajikan kepada publik.
\end{enumerate}

\subsection{Sistematika Penulisan}

Makalah ini disusun sebagai berikut. Bagian II membahas kajian
pustaka dan penelitian terkait mengenai sistem informasi geografis
untuk pengelolaan sampah perkotaan. Bagian III menjelaskan perancangan
sistem dan \textit{mockup} antarmuka pengguna. Bagian IV memaparkan
metode validasi, eksperimen, dan hasil yang diperoleh. Bagian V
menyajikan kesimpulan.

\section{Kajian Pustaka dan Penelitian Terkait}

\subsection{Penerapan GIS dalam Pengelolaan Sampah}

Penerapan Sistem Informasi Geografis (SIG/GIS) dalam pengelolaan
sampah perkotaan telah banyak dikaji dalam literatur internasional.
Asefa \textit{et al.} \cite{asefa2022} menguraikan bahwa GIS berperan
penting dalam menyederhanakan proses perencanaan pengelolaan sampah
yang kompleks, mulai dari optimasi rute pengumpulan hingga pemilihan
lokasi fasilitas, dengan menyediakan informasi spasial yang membantu
pengambil keputusan mengelola berbagai faktor yang saling bertentangan.
Pada tataran implementasi, Akpotuzor \textit{et al.} \cite{akpotuzor2020}
mengembangkan sistem pemantauan berbasis Web-GIS di Nigeria
menggunakan metodologi \textit{Y-model Web GIS Development}
(YWDM), yang memungkinkan warga melaporkan kondisi tempat sampah penuh
secara geospasial sekaligus mengoptimalkan rute truk pengangkut.
Pengujian fungsional (\textit{proof of concept}) pada sistem tersebut
menunjukkan bahwa pendekatan Web-GIS mampu menyediakan informasi
kondisi tempat sampah secara dinamis kepada pihak pengelola.

\subsection{Penelitian Terkait di Konteks Indonesia}

Pada konteks Indonesia, beberapa penelitian serupa telah dilakukan
dengan skala dan fokus yang berbeda-beda. Hidayat dan Nugraha
\cite{hidayat2023} mengembangkan sistem informasi manajemen pemetaan
bank sampah di Kota Batam untuk mengatasi ketiadaan sistem berbasis
peta yang membantu warga menemukan lokasi bank sampah terdekat.
Ardianti \textit{et al.} \cite{ardianti2025} mengembangkan aplikasi
pengelolaan sampah berbasis web terintegrasi GIS menggunakan metode
\textit{prototyping}, dengan fitur visualisasi peta interaktif
memakai Leaflet.js serta pelaporan transaksi --- pendekatan teknis
yang serupa dengan yang direncanakan pada penelitian ini. Pada domain
yang lebih spesifik, Fitriana \textit{et al.} \cite{fitriana2022}
mengembangkan aplikasi Android untuk studi kasus bank sampah Desa
Kalibagor, menunjukkan bahwa kebutuhan digitalisasi pengelolaan
sampah berbasis lokasi juga relevan pada skala pedesaan.

\subsection{Pendekatan Berbasis Sensor dan Partisipasi Warga}

Sebagai pembanding pendekatan teknis lain, Burman \textit{et al.}
\cite{burman2022} mengusulkan sistem manajemen sampah cerdas yang
memadukan GIS dengan sensor IoT pada tempat sampah untuk memantau
tingkat keterisian secara \textit{real-time} tanpa memerlukan input
manual. Pendekatan ini menawarkan presisi data yang tinggi namun
membutuhkan investasi perangkat keras yang signifikan pada setiap
titik pengumpulan. Di sisi lain, aplikasi seperti TrashOut
\cite{trashout} menunjukkan pendekatan \textit{crowdsourcing}, di
mana warga secara sukarela melaporkan lokasi dan kondisi sampah
melalui aplikasi mobile tanpa bergantung pada sensor maupun data
resmi pemerintah.

\subsection{Posisi Penelitian dan Kesenjangan}

Tabel \ref{tab:pustaka} merangkum perbandingan pendekatan pada
penelitian-penelitian terkait terhadap sistem yang diusulkan pada
makalah ini. Berdasarkan penelusuran pustaka, sebagian besar sistem
sejenis berfokus pada satu titik lokasi tempat sampah secara individual,
tanpa memodelkan secara eksplisit relasi satu fasilitas terhadap
banyak wilayah administratif (kelurahan) yang dilayaninya sekaligus
--- suatu karakteristik yang justru ditemukan secara konsisten pada
data internal TPS Kota Surabaya yang digunakan pada penelitian ini.
Selain itu, metode validasi pada penelitian terdahulu umumnya terbatas
pada pengujian fungsional atau iterasi \textit{prototyping}, tanpa
kombinasi validasi kualitas data, \textit{cross-check} terhadap basis
data eksternal resmi, pengujian usabilitas terstandardisasi (SUS), dan
\textit{User Acceptance Testing} (UAT) secara menyeluruh sebagaimana
dirancang pada Bagian IV makalah ini.

\begin{table}[h]
    \caption{Perbandingan dengan Penelitian Terkait}
    \label{tab:pustaka}
    \centering
    \begin{tabular}{p{1.5cm}p{1.4cm}p{2.1cm}p{2.2cm}}
        \hline
        \textbf{Penelitian}                            & \textbf{Platform} & \textbf{Sumber Data Kapasitas}                        & \textbf{Relasi Multi-Wilayah}                    \\
        \hline
        Akpotuzor \textit{et al.} \cite{akpotuzor2020} & Web-GIS           & Laporan warga manual                                  & Tidak eksplisit                                  \\
        \hline
        Hidayat \& Nugraha \cite{hidayat2023}          & Web-GIS           & Data lokasi statis                                    & Tidak ada                                        \\
        \hline
        Ardianti \textit{et al.} \cite{ardianti2025}   & Web + Leaflet.js  & Input petugas manual                                  & Tidak ada                                        \\
        \hline
        Burman \textit{et al.} \cite{burman2022}       & GIS + IoT         & Sensor \textit{real-time}                             & Tidak ada                                        \\
        \hline
        TrashOut \cite{trashout}                       & Mobile            & \textit{Crowdsourcing} warga                          & Tidak ada                                        \\
        \hline
        \textbf{Diusulkan}                             & Web-GIS           & Data internal terstruktur (231 entri) + laporan warga & \textbf{Eksplisit (many-to-many TPS--Kelurahan)} \\
        \hline
    \end{tabular}
\end{table}



\begin{thebibliography}{00}

    \bibitem{antara2023} Antara News Jawa Timur, ``DLH Surabaya catat 1.600 ton sampah di TPA Benowo didominasi jenis organik,'' 2023. [Online]. Tersedia: https://jatim.antaranews.com/berita/713517/

    \bibitem{liputan62023} Liputan6.com, ``Sampah di TPA Benowo Surabaya Capai 1.600 Ton per Hari, 60 Persen Didominasi Organik,'' 2023. [Online]. Tersedia: https://www.liputan6.com/surabaya/read/5338222/

    \bibitem{kompas2026} Kompas.com, ``Mengenal TPA Benowo yang Terbakar, Pusat Pengolahan Sampah Terbesar di Surabaya Sejak Tahun 2001,'' 2026. [Online]. Tersedia: https://surabaya.kompas.com/read/2026/07/20/

    \bibitem{suararakyat2026} Suara Rakyat Jatim, ``Komisi B Soroti Transisi Pengelolaan Sampah Surabaya, Pasca TPA Benowo Terbakar,'' 2026. [Online]. Tersedia: https://www.suararakyatjatim.com/2026/07/22/

    \bibitem{sipsn} Kementerian Lingkungan Hidup dan Kehutanan RI, ``Sistem Informasi Pengelolaan Sampah Nasional (SIPSN).'' [Online]. Tersedia: https://sipsn.menlhk.go.id

    \bibitem{worldbank} World Bank Group, ``What a Waste 2.0: A Global Snapshot of Solid Waste Management to 2050.'' [Online]. Tersedia: https://www.worldbank.org/en/publication/what-a-waste

    \bibitem{asefa2022} E. M. Asefa, K. B. Barasa, and D. A. Mengistu, ``Application of Geographic Information System in Solid Waste Management,'' in \textit{Geographic Information Systems and Applications in Coastal Studies}, IntechOpen, 2022. DOI: 10.5772/intechopen.103773.

    \bibitem{akpotuzor2020} S. A. Akpotuzor, A. O. Ofem, and M. A. Agana, ``Web-Based Monitoring System for Optimal Collection and Disposal of Solid Waste Using Geographic Information System (GIS),'' 2020. DOI: 10.5281/zenodo.3660159.

    \bibitem{hidayat2023} F. Hidayat and N. B. Nugraha, ``Optimizing the Waste Bank Mapping Management Information System in Batam City,'' \textit{Journal of Applied Geospatial Information}, vol. 7, no. 2, pp. 1080--1085, 2023. DOI: 10.30871/jagi.v7i2.4628.

    \bibitem{ardianti2025} M. Ardianti, S. Suakanto, and R. Z. Zulkarnaen, ``Development of Waste Management Application using GIS,'' \textit{SISTEMASI: Jurnal Sistem Informasi}, vol. 14, no. 6, pp. 3020--3027, 2025. DOI: 10.32520/stmsi.v14i6.5680.

    \bibitem{burman2022} A. R. Burman, A. Goswami, H. Sharma, and C. J. Mohan V, ``Smart Waste Management System using GIS Paired with IoT-based Smart Waste Collection Bins,'' in \textit{2022 Int. Conf. Smart and Sustainable Technologies in Energy and Power Sectors (SSTEPS)}, IEEE, 2022, pp. 229--232. DOI: 10.1109/SSTEPS57475.2022.00064.

    \bibitem{trashout} TrashOut, ``About TrashOut.'' [Online]. Tersedia: https://www.trashout.ngo/

    \bibitem{fitriana2022} G. F. Fitriana, A. Hashina, and N. A. F. Tanjung, ``Pengembangan Aplikasi Pengelolaan Sampah berbasis Android Studi Kasus Bank Sampah Desa Kalibagor,'' \textit{Journal of Dinda: Data Science, Information Technology, and Data Analytics}, vol. 2, no. 2, 2022. DOI: 10.20895/dinda.v2i2.741.

\end{thebibliography}

\end{document}